\documentclass[11pt]{article}

\usepackage[utf8]{inputenc}
\usepackage[T1]{fontenc}
\usepackage{amsmath, amssymb, amsthm}
\usepackage{mathtools}
\usepackage{geometry}
\usepackage{booktabs}
\usepackage{listings}
\usepackage{xcolor}
\usepackage{hyperref}
\usepackage{enumitem}

\geometry{margin=1in}

\definecolor{codebg}{rgb}{0.96,0.96,0.96}
\lstset{
  backgroundcolor=\color{codebg},
  basicstyle=\ttfamily\small,
  breaklines=true,
  frame=single,
  language=Python,
  showstringspaces=false,
  columns=fullflexible,
  keywordstyle=\color{blue!70!black}
}

\newcommand{\bp}{\mathbf{p}}
\newcommand{\bx}{\mathbf{x}}
\newcommand{\bw}{\mathbf{w}}
\newcommand{\bc}{\mathbf{c}}
\newcommand{\R}{\mathbb{R}}
\newcommand{\diver}{\nabla\!\cdot}
\newcommand{\grad}{\nabla}

\title{RBF Finite-Difference Discretization of Anisotropic Diffusion:\\
Theory and Implementation}
\author{Dennis Corraliza}
\date{6/21/2026}

\begin{document}
\maketitle

\begin{abstract}
This note derives the Radial Basis Function Finite Difference (RBF-FD)
method for the anisotropic diffusion equation
$-\diver(A\grad u) = f$ on a bounded domain $\Omega \subset \R^2$,
with mixed Dirichlet/Neumann boundary conditions, and shows precisely how
each step of the derivation corresponds to a function in the
accompanying Python implementation (modules \texttt{assembly.py},
\texttt{domain.py}, \texttt{rbf.py}, \texttt{boundary.py},
\texttt{stencils.py}, and \texttt{geometry.py}).
\end{abstract}

\tableofcontents

%==============================================================
\section{Problem statement}
%==============================================================

We seek $u : \Omega \to \R$ satisfying the anisotropic diffusion
equation
\begin{equation}
  -\diver\!\big(A \grad u\big) = f \quad \text{in } \Omega,
  \label{eq:pde}
\end{equation}
where $A \in \R^{2\times 2}$ is a constant, symmetric positive-definite
diffusion tensor and $f$ is a prescribed source term. The boundary
$\partial\Omega$ is partitioned into Dirichlet and Neumann pieces:
\begin{align}
  u &= g_D \quad \text{on } \Gamma_D \subset \partial\Omega, \\
  \frac{\partial u}{\partial n} \equiv n \cdot \grad u &= g_N \quad
  \text{on } \Gamma_N \subset \partial\Omega,
  \label{eq:neumann-bc}
\end{align}
where $n$ is the outward unit normal. When $A = I$, \eqref{eq:pde}
reduces to the standard Poisson equation $-\Delta u = f$.

In the code, the geometric domain is either a square $[0,L]^2$ or a
disk of radius $R$ (\texttt{geometry.py}), the tensor $A$ is supplied
directly by the user, and the boundary type on each side/arc is
encoded by the list \texttt{btype} (e.g.\ \texttt{['dirichlet',
'neumann', 'dirichlet', 'neumann']} for a square, ordered
counterclockwise starting at $y=0$).

%==============================================================
\section{Radial basis function interpolation}
%==============================================================

\subsection{Local RBF interpolant}

Let $\{\bx_j\}_{j=1}^{N}$ be the full node set discretizing $\Omega$.
For a fixed node $\bx_i$, let $S_i = \{j_1,\dots,j_n\} \subset
\{1,\dots,N\}$ index its \emph{stencil}: the $n$ nearest neighbors of
$\bx_i$ (including $\bx_i$ itself), determined by a $k$-nearest-neighbor
search. We build a local interpolant of a function $u$ over the
stencil using a radial kernel $\varphi$:
\begin{equation}
  u(\bx) \;\approx\; \sum_{j \in S_i} \lambda_j\, \varphi(\|\bx - \bx_j\|)
  \;+\; \sum_{\ell=1}^{m} \mu_\ell\, p_\ell(\bx),
  \label{eq:interp}
\end{equation}
where $\{p_\ell\}_{\ell=1}^m$ is an (optional) polynomial basis used to
\emph{augment} the kernel interpolant, with coefficients constrained by
\begin{equation}
  \sum_{j \in S_i} \lambda_j\, p_\ell(\bx_j) = 0, \qquad \ell = 1,\dots,m.
  \label{eq:augment-constraint}
\end{equation}
Two kernels are used in the implementation (\texttt{rbf.py}):
\begin{align}
  \varphi_{\text{cubic}}(\bp) &= \|\bp\|^3,
  &\varphi_{\text{gauss}}(\bp) &= \exp\!\big(-(\varepsilon\|\bp\|)^2\big),
  \label{eq:kernels}
\end{align}
where $\bp = \bx - \bx_j$ is a displacement vector and $\varepsilon$ is
a user-chosen shape parameter for the Gaussian. The polynomial basis
used for augmentation is linear,
\begin{equation}
  \{p_1, p_2, p_3\}(x,y) = \{1,\ x,\ y\},
  \label{eq:polybasis}
\end{equation}
so that the interpolant exactly reproduces constants and linear
functions whenever augmentation is enabled.

\paragraph{Code correspondence.}
\begin{itemize}[leftmargin=2em]
  \item \texttt{rbf.phi\_cubic(p)}, \texttt{rbf.phi\_gauss(p, eps)}
        $\;\leftrightarrow\;$ \eqref{eq:kernels}.
  \item \texttt{rbf.poly\_basis(p)} $\;\leftrightarrow\;$
        \eqref{eq:polybasis}.
  \item \texttt{stencils.knn\_list(P, num\_stencil\_nodes)} builds $S_i$ for every node.
  \item \texttt{domain.PDEDomainContext.augmentation} is the boolean
        flag toggling the second sum in \eqref{eq:interp}.
\end{itemize}

%==============================================================
\section{RBF-FD differentiation weights}
%==============================================================

\subsection{The linear functional viewpoint}

RBF-FD does not differentiate the interpolant \eqref{eq:interp}
symbolically; instead, it seeks weights $w_j$ such that a linear
operator $\mathcal{L}$ applied at $\bx_i$ is approximated directly as a
weighted sum of \emph{nodal values}:
\begin{equation}
  (\mathcal{L}u)(\bx_i) \;\approx\; \sum_{j \in S_i} w_j\, u(\bx_j).
  \label{eq:weighted-sum}
\end{equation}
The weights are chosen so that \eqref{eq:weighted-sum} is \emph{exact}
whenever $u$ is any kernel function centered at a stencil node, and
(if augmentation is used) exact for every polynomial in the
augmentation basis. This yields the linear system
\begin{equation}
  \underbrace{
  \begin{bmatrix} M & P \\ P^\top & 0 \end{bmatrix}
  }_{\text{generalized Gram matrix}}
  \begin{bmatrix} \bw \\ \boldsymbol{\mu} \end{bmatrix}
  =
  \begin{bmatrix} \mathbf{b} \\ \mathbf{0} \end{bmatrix},
  \label{eq:local-system}
\end{equation}
where
\begin{align}
  M_{jk} &= \varphi(\bx_{j} - \bx_{k}), &j,k \in S_i, \\
  P_{j\ell} &= p_\ell(\bx_j), &j \in S_i,\ \ell=1,\dots,m, \\
  b_j &= (\mathcal{L}\varphi)(\bx_i - \bx_j), &j \in S_i.
  \label{eq:local-system-entries}
\end{align}
The vector $\mathbf{b}$ holds the operator $\mathcal{L}$ applied
\emph{exactly} (in closed form) to the kernel, evaluated at the
displacement from the stencil center to each stencil node, rather than
any discrete approximation. Solving \eqref{eq:local-system} for $\bw$
gives the desired differentiation weights at node $\bx_i$.

\paragraph{Code correspondence.}
This is exactly what \texttt{assembly.local\_weights\_solve(context,i)}
and \texttt{assembly.local\_grad\_solve(context,i)} implement, for
$\mathcal{L} = \diver(A\grad\,\cdot\,)$ and $\mathcal{L} = \grad$
respectively:
\begin{itemize}[leftmargin=2em]
  \item The matrix \texttt{M} in the code is exactly $M$ in
        \eqref{eq:local-system-entries}, built from
        \texttt{context.phi}.
  \item The vector \texttt{b} (or matrix \texttt{b\_grad}, with one
        column per spatial dimension) is built from \\
        \texttt{context.laplacian\_phi} (resp.\ \texttt{context.grad\_phi}),
        which are the closed-form expressions $\diver(A\grad\varphi)$
        and $\grad\varphi$ derived in Section~\ref{sec:aniso-op} below.
  \item The block matrix \texttt{Pmat} and the block-augmented system
        built via \texttt{np.block} are exactly the augmented system
        in \eqref{eq:local-system} when \texttt{context.augmentation}
        is \texttt{True}.
  \item \texttt{np.linalg.solve(M, b)} solves \eqref{eq:local-system}
        and the function returns only the first \texttt{num\_nodes}
        entries of the solution, i.e.\ $\bw$ (discarding the
        polynomial Lagrange multipliers $\boldsymbol{\mu}$).
\end{itemize}

\subsection{Least-squares variant}
\label{sec:ls-variant}

An alternative formulation avoids solving the square system
\eqref{eq:local-system} directly and instead seeks weights that best
reproduce $\mathcal{L}$ in a least-squares sense over a richer set of
\emph{auxiliary evaluation centers} $\{\bc_k\}_{k=1}^{K}$ scattered
around $\bx_i$ (typically more centers than stencil nodes, $K > n$):
\begin{equation}
  \min_{\bw}
  \left\|
    \underbrace{\big[\varphi(\bc_k - \bx_j)\big]_{k,j}}_{=:M \in
    \R^{K\times n}} \bw
    \;-\;
    \underbrace{\big[(\mathcal{L}\varphi)(\bx_i - \bc_k)\big]_k}_{=:\mathbf{b}
    \in \R^{K}}
  \right\|_2^2,
  \label{eq:ls-system}
\end{equation}
subject to the same polynomial reproduction constraints
\eqref{eq:augment-constraint} enforced exactly (as equality rows
appended to the system, rather than penalized in the least-squares
objective). The motivation for this variant is robustness: anchoring
the fit to a denser, more uniformly distributed cloud of auxiliary
points reduces sensitivity to clustering or anisotropy in the stencil
node positions themselves.

The auxiliary centers $\{\bc_k\}$ are generated as a near-uniform
polar grid of $num_rings$ concentric rings (spacing $h = r/num_rings$, with $r$ the
stencil radius) plus a center point, with the number of angular
samples per ring chosen so that the arc-length spacing on each ring is
approximately $h$ as well:
\begin{equation}
  \bc_k = \big(\rho_i \cos\theta_{ij},\ \rho_i \sin\theta_{ij}\big),
  \qquad \rho_i = ih,\ \ \theta_{ij} = \frac{2\pi j}{n_i},
  \quad n_i = \max\!\Big(1,\ \mathrm{round}\big(\tfrac{2\pi\rho_i}{h}\big)\Big).
  \label{eq:aux-grid}
\end{equation}

\paragraph{Code correspondence.}
\begin{itemize}[leftmargin=2em]
  \item \texttt{rbf.generate\_grid\_2d(r, k)} implements
        \eqref{eq:aux-grid}, using \texttt{rbf.rad\_to\_euc} to convert
        each $(\rho_i, \theta_{ij})$ to Cartesian coordinates.
  \item \texttt{assembly.local\_weights\_ls(context,i)} and
        \texttt{assembly.local\_grad\_ls(context,i)} \\ build $M$ and
        $\mathbf{b}$ exactly as in \eqref{eq:ls-system}, append the
        polynomial constraint rows, and call \\
        \texttt{np.linalg.lstsq} to solve.
  \item Dispatch between the direct system \eqref{eq:local-system} and
        the least-squares system \eqref{eq:ls-system} is controlled by
        \texttt{context.center\_rings}: \texttt{None} selects the
        direct solve, any integer value selects the least-squares
        variant with that many rings $num_rings$.
\end{itemize}

%==============================================================
\section{The anisotropic diffusion operator applied to each kernel}
\label{sec:aniso-op}
%==============================================================

The entries $b_j$ in \eqref{eq:local-system-entries} require a
closed-form expression for $\diver(A\grad\varphi)$, evaluated at the
displacement $\bp = \bx_i - \bx_j$. Writing $A = \begin{psmallmatrix}
A_{11} & A_{12} \\ A_{12} & A_{22} \end{psmallmatrix}$ and
$r = \|\bp\|$, $\bp = (x,y)$:

\paragraph{Cubic kernel.} For $\varphi(\bp) = r^3$,
\begin{equation}
  \diver\!\big(A\grad\varphi\big)(\bp)
  = 3r\left(\operatorname{tr}(A) + \frac{\bp^\top A \bp}{r^2}\right).
  \label{eq:lap-cubic}
\end{equation}
This expression is regular as $r \to 0$ in the sense that the
true value of the operator at $\bp=0$ is $0$; the implementation
returns $0$ directly whenever $r$ falls below a tolerance, avoiding
the removable $1/r^2$ singularity in \eqref{eq:lap-cubic} without
altering the limiting value.

\paragraph{Gaussian kernel.} For $\varphi(\bp) =
\exp(-(\varepsilon r)^2)$,
\begin{equation}
  \diver\!\big(A\grad\varphi\big)(\bp)
  = 2\varepsilon^2\,\varphi(\bp)
  \Big(2\varepsilon^2\, \bp^\top A \bp - \operatorname{tr}(A)\Big).
  \label{eq:lap-gauss}
\end{equation}
No singularity occurs here since $\varphi$ and its derivatives are
smooth everywhere.

\paragraph{Polynomial augmentation basis.} Since
$\diver(A\grad p_\ell) \equiv 0$ for $p_\ell \in \{1,x,y\}$ (all
second derivatives vanish), the corresponding entries in the augmented
right-hand side of \eqref{eq:local-system} are identically zero,
independent of $A$.

\paragraph{Gradients.} The two kernel gradients used to build Neumann
boundary rows (Section~\ref{sec:neumann}) are
\begin{equation}
  \grad\varphi_{\text{cubic}}(\bp) = 3r\,\bp,
  \qquad
  \grad\varphi_{\text{gauss}}(\bp) = -2\varepsilon^2\,\varphi_{\text{gauss}}(\bp)\,\bp.
  \label{eq:grad-kernels}
\end{equation}

\paragraph{Code correspondence.}
\begin{itemize}[leftmargin=2em]
  \item \texttt{rbf.anisotropic\_diffusion\_phi\_cubic(p, A, tol)}
        $\leftrightarrow$ \eqref{eq:lap-cubic}; \\
        \texttt{rbf.anisotropic\_diffusion\_phi\_gauss(p, A, eps)}
        $\leftrightarrow$ \eqref{eq:lap-gauss}.
  \item \texttt{rbf.anisotropic\_diffusion\_poly(p, A, tol)} returns
        \texttt{0.0}, matching the polynomial-basis observation above.
  \item \texttt{rbf.grad\_phi\_cubic(p)}, \texttt{rbf.grad\_phi\_gauss(p, eps)}
        $\leftrightarrow$ \eqref{eq:grad-kernels}.
  \item \texttt{assembly.set\_rbf\_func} wires the chosen kernel,
        gradient, and anisotropic-diffusion closures (with $A$ fixed)
        onto \texttt{context.phi}, \texttt{context.grad\_phi}, and
        \texttt{context.laplacian\_phi} respectively, dispatching on
        the string \texttt{basis} (\texttt{'cubic'} or
        \texttt{'gaussian'}).
\end{itemize}

%==============================================================
\section{Global assembly}
%==============================================================

\subsection{Interior rows}

For every node $\bx_i$ in the domain, the local weight vector $\bw^{(i)}
\in \R^{|S_i|}$ computed by solving \eqref{eq:local-system} (or
\eqref{eq:ls-system}) is scattered into row $i$ of a dense global
matrix $W \in \R^{N \times N}$:
\begin{equation}
  W_{ij} =
  \begin{cases}
    w^{(i)}_{j} & j \in S_i, \\
    0 & j \notin S_i.
  \end{cases}
  \label{eq:global-W}
\end{equation}
At this stage, $W$ approximates $\diver(A\grad\,\cdot\,)$ at every node,
ignoring boundary conditions entirely.

\paragraph{Code correspondence.}
\texttt{assembly.global\_weights(context, in\_boundary, normal\_vec)} \\
loops over all stencils, dispatches to
\texttt{local\_weights\_solve} or \texttt{local\_weights\_ls} depending
on \\ \texttt{context.center\_rings}, and assembles \eqref{eq:global-W}.

\subsection{Dirichlet rows}

For a node $\bx_i \in \Gamma_D$, row $i$ of $W$ is overwritten with the
identity row,
\begin{equation}
  W_{ij} = \delta_{ij}, \qquad i \in \Gamma_D,
  \label{eq:dirichlet-row}
\end{equation}
so that the linear system enforces $u(\bx_i) = g_D(\bx_i)$ exactly
(via the matching right-hand-side entry, Section~\ref{sec:rhs}).

\subsection{Neumann rows}
\label{sec:neumann}

For a node $\bx_i \in \Gamma_N$ with outward unit normal $n_i$, row $i$
is overwritten with directional-derivative weights, built from the
\emph{local gradient} weight matrix $W^{(i)}_{\text{grad}} \in
\R^{|S_i| \times 2}$ (the analog of \eqref{eq:local-system} for
$\mathcal{L} = \grad$):
\begin{equation}
  W_{ij} = \Big(W^{(i)}_{\text{grad}}\, n_i\Big)_{j}, \qquad j \in S_i,
  \quad i \in \Gamma_N,
  \label{eq:neumann-row}
\end{equation}
so that row $i$ of the assembled system enforces $n_i \cdot \grad u
(\bx_i) \approx g_N(\bx_i)$, matching \eqref{eq:neumann-bc}.

\paragraph{Code correspondence.}
\texttt{assembly.boundary\_to\_weights} implements \eqref{eq:dirichlet-row} and
\eqref{eq:neumann-row} in place, dispatching per node on the label
returned by \texttt{in\_boundary(p)} (itself built in \\
\texttt{assembly.set\_boundary\_func} from
\texttt{boundary.in\_square\_boundary} / \\
\texttt{boundary.in\_circular\_boundary}, with normals from
\texttt{boundary.square\_normal} / \\ \texttt{boundary.disk\_normal}).

\subsection{Right-hand side}
\label{sec:rhs}

The right-hand-side vector $\mathbf{F} \in \R^N$ is assembled
nodewise,
\begin{equation}
  F_i =
  \begin{cases}
    f(\bx_i) & \bx_i \text{ interior}, \\
    g_D(\bx_i) & \bx_i \in \Gamma_D, \\
    g_N(\bx_i) & \bx_i \in \Gamma_N,
  \end{cases}
  \label{eq:rhs}
\end{equation}
so that the assembled linear system $W\mathbf{u} = \mathbf{F}$ is
exactly the discrete analog of \eqref{eq:pde}--\eqref{eq:neumann-bc}.

\paragraph{Code correspondence.}
\texttt{assembly.right\_hand\_side(context,f,g,in\_boundary)}
implements \eqref{eq:rhs}, where the single callable \texttt{g} passed
in already dispatches internally to the correct side/arc function
(see \texttt{boundary.square\_boundary}, \texttt{boundary.circular\_boundary}).

%==============================================================
\section{The pure-Neumann null space and anchoring}
%==============================================================

If $\Gamma_D = \emptyset$ (every boundary node is Neumann), the
continuous problem \eqref{eq:pde}--\eqref{eq:neumann-bc} determines $u$
only up to an additive constant, since $\diver(A\grad c) = 0$ and
$n\cdot\grad c = 0$ for any constant $c$. The discrete matrix $W$
inherits this null space and is singular. Three regularization
strategies are implemented:
\begin{align}
  \textbf{mean: } &\textstyle\sum_i u_i = 0,
  &&\text{(replace the last row of $W$ by $\mathbf{1}^\top$, set
  $F_{\text{last}}=0$),} \\
  \textbf{pin: } &u_0 = 0,
  &&\text{(replace the first row of $W$ by $\mathbf{e}_0^\top$, set
  $F_0 = 0$),} \\
  \textbf{project: } &W \leftarrow \Pi W \Pi,\ \ \mathbf{F} \leftarrow
  \Pi \mathbf{F},
  &&\Pi = I - \tfrac1N \mathbf{1}\mathbf{1}^\top.
  \label{eq:anchor}
\end{align}

\paragraph{Code correspondence.}
\texttt{assembly.is\_pure\_neumann(context,in\_boundary)} detects this
case; \texttt{assembly.anchor\_system(W, f, method)} implements the
three strategies in \eqref{eq:anchor}, selected by the string
\texttt{method} (\texttt{'mean'}, \texttt{'pin'}, or \texttt{'project'}).

%==============================================================
\section{Solving and evaluating the discrete solution}
%==============================================================

The fully assembled (and, if necessary, anchored) linear system
\begin{equation}
  W\mathbf{u} = \mathbf{F}
  \label{eq:final-system}
\end{equation}
is solved directly for the nodal solution vector $\mathbf{u} \in
\R^N$, approximating $u(\bx_i)$ at every node.

\paragraph{Code correspondence.}
\texttt{assembly.rbf\_fd\_solve(W, F)} calls
\texttt{np.linalg.solve(W, F)}, implementing \eqref{eq:final-system}.
The full pipeline --- nodes and stencils, kernel and boundary
configuration, weight assembly, boundary imposition, right-hand-side
assembly, and anchoring --- is orchestrated end-to-end by
\texttt{assembly.rbf\_fd\_system(\dots)}, which returns a populated
\texttt{domain.PDEDomainContext} carrying $A$, $W$, and $\mathbf{F}$
for later use by \texttt{rbf\_fd\_solve}.

%==============================================================
\section{Convergence diagnostics}
%==============================================================

For manufactured-solution problems (Section~\ref{sec:manufactured}),
the discrete error is reported in both the maximum norm and a
discrete $L^2$ norm:
\begin{equation}
  \|e\|_\infty = \max_i |u_i - u_{\text{exact}}(\bx_i)|,
  \qquad
  \|e\|_{2} = \frac{1}{\sqrt{N}}
  \left(\sum_{i=1}^N \big(u_i - u_{\text{exact}}(\bx_i)\big)^2\right)^{1/2},
  \label{eq:errors}
\end{equation}
together with the condition number $\kappa(W) =
\|W\|\,\|W^{-1}\|$ as a diagnostic for ill-conditioning of the
assembled system (which tends to grow as stencils shrink or as $num_rings$,
$\varepsilon$, or node spacing are pushed to extremes).

\paragraph{Code correspondence.}
These are exactly the quantities computed and printed in \\
\texttt{report\_and\_graph(context, u\_exact)}: \texttt{np.linalg.cond(W)}
for $\kappa(W)$, \texttt{np.max(error)} for $\|e\|_\infty$, and
\texttt{np.linalg.norm(error)/np.sqrt(len(P))} for $\|e\|_2$, alongside
contour and surface plots of the approximate solution, exact solution,
and pointwise error.

%==============================================================
\section{Manufactured solutions}
\label{sec:manufactured}
%==============================================================

To validate the method, exact solutions $u_{\text{exact}}$ are chosen
in closed form, and the corresponding forcing term $f$ and boundary
data $g_D, g_N$ are derived analytically so that $u_{\text{exact}}$
satisfies \eqref{eq:pde}--\eqref{eq:neumann-bc} exactly (the
\emph{method of manufactured solutions}). For example, the
mixed polynomial--trigonometric solution
\begin{equation}
  u_{\text{exact}}(x,y) =
  (L-x)\Big(y - \tfrac{y^2}{L}\Big)
  + c\,(L-y)\,\frac{x(L-x)}{L}
  + \mathrm{Amp}\,\sin\!\Big(\frac{\alpha\pi x}{L}\Big)
  \sin\!\Big(\frac{\beta\pi y}{L}\Big),
  \quad c=2,
  \label{eq:manufactured-u}
\end{equation}
yields, after applying $-\diver(A\grad\,\cdot\,)$ to
\eqref{eq:manufactured-u} and collecting polynomial, trigonometric, and
mixed-derivative contributions, a forcing term $f$ with three additive
pieces:
\begin{align}
  f_{\text{poly}} &= -2cA_{11}\frac{L-y}{L}
  + 2A_{12}\Big(\frac{2cx}{L} + \frac{2y}{L} - \frac{c+1}{L}\Big)
  - 2A_{22}\frac{L-x}{L}, \\
  f_{\text{trig}} &= -\mathrm{Amp}\Big(A_{11}k_x^2 + A_{22}k_y^2\Big)
  \sin(k_x x)\sin(k_y y), \\
  f_{\text{mixed}} &= 2\,\mathrm{Amp}\,A_{12}\,k_x k_y\,\cos(k_x x)\cos(k_y y),
  \label{eq:manufactured-f}
\end{align}
with $k_x = \alpha\pi/L$, $k_y = \beta\pi/L$, and $f = f_{\text{poly}} +
f_{\text{trig}} + f_{\text{mixed}}$, together with Dirichlet data read
off directly from \eqref{eq:manufactured-u} on each edge.

\paragraph{Code correspondence.}
\texttt{example\_3(Amp,modes,A,L)} returns the tuple
\texttt{(f,g,btype,u\_exact)} implementing exactly
\eqref{eq:manufactured-u}--\eqref{eq:manufactured-f}: \texttt{u\_exact}
evaluates \eqref{eq:manufactured-u}, \texttt{f} evaluates the sum of
\eqref{eq:manufactured-f}, and \texttt{g} is the list of four
one-variable boundary functions obtained by restricting
\eqref{eq:manufactured-u} to each edge of the square.

%==============================================================
\section{Node generation strategies}
%==============================================================

The implementation supports several node distributions
(\texttt{geometry.py}):
\begin{itemize}[leftmargin=2em]
  \item \textbf{Uniform tensor-product grid} (\texttt{uniform\_square}):
        a regular $N_x \times N_y$ Cartesian grid over $[0,L]^2$,
        boundary and interior nodes undifferentiated.
  \item \textbf{Uniform interior + explicit boundary}
        (\texttt{uniform\_int\_square}): interior nodes from a regular
        grid strictly inside $[0,L]^2$, concatenated with $N_b$
        equally spaced nodes per edge (corners de-duplicated); returns
        the interior node count alongside the array so that
        boundary/interior can be distinguished downstream if needed.
  \item \textbf{Chebyshev--Gauss--Lobatto grid} (\texttt{cheby\_square}):
        tensor-product nodes \\ $x_j = \tfrac{L}{2}\big(1+\cos(j\pi/N_x)\big)$,
        clustering near the edges, as is standard for spectral
        collocation methods.
  \item \textbf{Circular domain} (\texttt{circular\_geometry}):
        interior nodes from a Cartesian grid filtered to the open disk
        $x^2+y^2<R^2$, plus boundary nodes placed exactly on the circle
        at $N_{\text{bound}}$ equally spaced angles.
\end{itemize}
In every case, the resulting node array $P \in \R^{N\times 2}$ is the
sole input to \texttt{domain.PDEDomainContext}, after which stencils
(\texttt{stencils.knn\_list}) and RBF-FD weights are computed purely
from the node positions, independent of how they were generated.

\end{document}
